С помощью различных датчиков РНТТМ обеспечивается сбор данных о состоянии окружающей среды. Датчики можно классифицировать на датчики внешней информации и внутренней \cite{Vlasov}.

Датчики внутренней информации РНТТМ представляют собой в основном преобразователи механических параметров в электрические сигналы. Например датчики обратной связи по положению, скорости и ускорения.

\begin{figure}[ht]
    \centering{
        \includegraphics[scale=0.45]{drawio/datchiki.drawio.png}
    }
    \caption{Классификация датчиков внешней информации РНТТМ}\label{fig:datchiki}
\end{figure}

Датчики внешней информации РНТТМ обеспечивают сбор информации об окружающем мире с различными целями. Датчики можно разделить на две большие группы (рис.\ \cref{fig:datchiki}): дистанционные и контактные. В свою очередь дистанционные датчики можно разделить на локационные и визуальные.

Локационные системы РНТТМ можно условно разделить на два класса: дальней и~ближней локации. Первые могут быть построены с использованием ультразвуковых, лазерных и светолокационных датчиков.

Визуальные системы (технического, машинного зрения) РНТТМ относятся к числу наиболее сложных, но и наиболее универсальных средств очувствления РНТТМ (рис.\ \cref{fig:ypiter30},\ \cref{fig:robot1},\ \cref{fig:robot2}).

\begin{figure}[ht]
	\centering{
		\hfill
		\subcaptionbox{Робота\label{fig:robot1}}{%
			\includegraphics[width=0.4\linewidth]{views/robot1.jpg}}
		\hfill
		\subcaptionbox{Кобота\label{fig:robot2}}{%
			\includegraphics[width=0.4\linewidth]{views/robot2.jpg}}
		\hfill
	}
	\caption{Примеры наборов датчиков}\label{fig:robot1_2}
\end{figure}

Основными используемыми средствами машинного видения РНТТМ~\cite{confbib1} являются:

Лазерные дальномеры (лидары) РНТТМ, которые измеряют расстояние до объектов. Результат работы лидара -- объемная сцена из облака точек с геометрией объектов. Это средство выдает наибольшую точность, лучше всех работают на ровных поверхностях и почти не засвечиваются солнцем. Недостатки лидаров -- это высокое энергопотребление, относительно низкая частота кадров, бегущий затвор и необходимость компенсировать его при обработке, а также работающие рядом лидары создают друг-другу помехи, которые не так просто компенсировать.

Радары РНТТМ для определения расстояния до объекта, его скорости и~местораспо-ложения. Радары меньше зависят от погоды и стоят намного дешевле лидаров. Главным отрицательным качеством радаров является то, что плохо обнаруживаются неметаллические объекты, в том числе пешеходы.

Инфракрасные датчики РНТТМ для реагирования на фоновое инфракрасное излучение. Прибор регистрирует любое тепловое излучение.

Time-of-flight-камеры (ToF) РНТТМ -- видеокамеры, формирующие дальностное изображение. Используются для создания изображений, которые в качестве пикселей содержат оценки расстояний от экрана до конкретных точек наблюдения.

\nomenclature{ToF}{Time-of-flight-камеры\nomrefpage}

Стерео-камеры РНТТМ, построенные на имитации бинокулярного зрения человека с возможностью захватывать трехмерные изображения.

Пленоптические камеры РНТТМ, фиксирующие векторное поле световых лучей (световое поле). На основе картины светового поля может быть воссоздана наиболее полная информация об изображении, пригодная для решения различных задач компьютерной графики.

Обычные камеры РНТТМ с последующим анализом видео-потока с применением алгоритмов машинного обучения.

Применение какого-либо только одного средства недостаточно, так как получение комплексного видения окружающего мира роботами имеет сложности в виде:

\begin{itemize}[beginpenalty=10000, midpenalty=10000, wide=\parindent, leftmargin=0pt]
    \item темного времени суток;
    \item засветов от солнца и источников освещения;
    \item инфракрасного излучение солнечного света;
    \item осадков и грязи;
    \item множественности и вариативности подвижных окружающих объектов.
\end{itemize}

В связи с этим, средства машинного видения в РНТТМ применяются комплексно, для компенсации перечисленных недостатков. Сравнительный анализ приведен в~таблице~\ref{tab:TableViews}. Согласно требованиям к подсистеме распознавания ГОСТ Р 71615-2024 (п. 6.4.10) \cite{gostR71615}, для обеспечения функциональной безопасности ВАЗДСМ необходимо комбинирование указанных средств технического зрения (лидары, радары, камеры) для исключения ложных срабатываний в условиях запыленности и осадков.

\begin{table}[htbp]
	\centering
	\caption{Сравнительный анализ средств машинного видения РНТТМ}\label{tab:TableViews}
	\resizebox{\textwidth}{!}{%
		\begin{threeparttable}
			\begin{tabular}{|l|c|c|c|c|c|}
				\hline
				& \textbf{ToF} & \textbf{Стерео} & \textbf{Пленопт.} & \textbf{Лидары} & \textbf{Радары} \\ \hline
				Разрешение                 & \textit{Слабо}  & Средне & \textbf{Отлично (сцена)} & Средне (сцена) & \textit{Слабо}  \\ \hline
				Точность                  & Средне & Средне & \textit{Слабо} & \textit{Слабо} & \textbf{Высоко} \\ \hline
				Сложность программ        & \textbf{Слабо} & Средне & \textit{Высоко} & \textit{Высоко} & \textbf{Слабо} \\ \hline
				Работа в реальном времени & \textbf{Высоко} & Средне & Средне & Средне & \textit{Слабо} \\ \hline
				Стоимость                 & Средне & Средне & \textbf{Слабо} & \textit{Высоко} & Средне \\ \hline
				Произв. при слабом освещении & \textbf{Хорошо} & \textbf{Хорошо} & \textit{Плохо} & \textit{Плохо} & \textbf{Хорошо} \\ \hline
				Работа на открытом воздухе & \textit{Плохо} & \textit{Плохо} & \textbf{Хорошо} & \textbf{Хорошо} & \textbf{Хорошо} \\ \hline
			\end{tabular}
		\end{threeparttable}
	}
\end{table}

Можно сделать вывод, что по разрешению лидируют пленоптические камеры, но все очень сильно зависит от состояния окружающей среды (сцены съемки). По точности -- лидары в РНТТМ. По сложности обработки -- «непосредственно» получают глубину только ToF и~лидары. По частоте кадров конструктивно хороши ToF-камеры. При работе на открытом пространстве плохо работают ToF и~камеры структурированного света.